%% notes-tex/database/pputida-cell-state/sections/isoprenol.tex
%% [[experiments.database.pputida-cell-state]]
%% https://github.com/Mjvolk3/torchcell/tree/main/notes-tex/database/pputida-cell-state/sections/isoprenol.tex
%% SOURCE: experiments/database/scripts/build_pputida_cell_state_table.py
%% (Tables pcsiso, pcscycle and pcsrawcycle); the prose is hand-authored.

\section{Predicting the last Carruthers round\secstatus{todo}}\label{sec:isoprenol}

The task these rows were surveyed for: predict the final design-build-test-learn round
of the Carruthers isoprenol campaign from its earlier rounds, using whatever other
isoprenol data is served. Carruthers is a collaborator, so the campaign's own data is
available; the question is which single dataset from outside it helps.

\subsection{The dataset: Yunus 2026\secstatus{todo}}\label{sec:iso-yunus}

\texttt{IsoprenolTiterYunus2026} is the answer, and Table~\ref{tab:pcsiso} is why.
Yunus 2026 (doi:\allowbreak 10.1016/\allowbreak j.ymben.\allowbreak 2025.11.007)
selects CRISPRi knockdown targets by FluxRETAP against intuition, in the engineered
isoprenol-producing strain IY1452, and Carruthers is a co-author. Four things line up
and one does not.

\notesec{The target sets nearly coincide} Every target in the Yunus campaign but one
is also a Carruthers target (Table~\ref{tab:pcsiso}). That is what makes the dataset
unusual: it is not a different experiment that happens to share a product, it is a
second measurement of nearly the same target set.

\notesec{The readout is the same quantity} Same product, same unit, same
quantification method, same 48 hour culture. No unit conversion and no method bridge
stands between the two campaigns. The replicate structure is the one readout axis that
differs: Carruthers varies from 2 to 6 biological replicates per strain and every
record carries a sample standard deviation, while Yunus is uniformly 3 and no record
carries a dispersion at all (Table~\ref{tab:pcsiso}).

\notesec{The chassis and the modality are the same} Both repress a native gene by
CRISPRi in an isoprenol-producing strain, so a target's effect means the same thing in
both.

\notesec{The paired features exist} \sym{CrispriPanelProteomeYunus2026} carries an
absolute Top3 proteome over 253 accessions on the same strains, which is the same
feature modality as \sym{CampaignProteomeCarruthers2025} in
Table~\ref{tab:pcsrows}. A proteome-to-titer map can be fit on both.

\notesec{The conditions differ, and this decides how it is used} Carruthers runs at
24\,\textdegree C in M9-NREL with 20 g/L glucose; Yunus runs at 30\,\textdegree C in
plain M9 salts. The titer ranges differ accordingly, and the Yunus records carry no
uncertainty, because only one number per strain was released and the loader records
that as a typed gap rather than inventing a dispersion.

So Yunus is not extra training rows for absolute titer. Two uses survive the condition
difference, and both are worth more than 125 rows would be:

\begin{itemize}[nosep,leftmargin=1.4em]
\item \textbf{A held-out condition.} If a model trained on the earlier Carruthers
cycles ranks targets correctly in Yunus, the ranking is a property of the biology
rather than of the campaign. Test rank agreement, not absolute error, since the scales
differ.
\item \textbf{A reproducibility bound on the label.} Two independent measurements of
nearly the same target set say how much of the titer variation is real. That ceilings
what any model can achieve on this data, and it is the first thing to compute.
\end{itemize}

The runner-up is \sym{IsoprenolSelectionMenasalvas2025} in
Table~\ref{tab:pcsrows}, the only other CRISPRi and isoprenol pairing. Its readout is a
categorical biosensor enrichment call rather than a titer, so it is a weaker label and
not the first choice.

\subsection{The blocker: the cycle label is not in the graph\secstatus{todo}}\label{sec:iso-cycle}

The task needs a cycle label, to train on the earlier rounds and test on the last one.
No served Carruthers record carries one. Table~\ref{tab:pcscycle} searches every record
in all four Carruthers stores for a DBTL token anywhere in its serialization and finds
none, and \texttt{screen\_id} is empty throughout.

The loader's record identity is one per construct and DBTL cycle, so the cycle is
implied by the record set and was never written onto a record. From the graph alone the
final round cannot be separated from the earlier ones, which means the task as stated
cannot be set up from a query.

It is recoverable. The raw mirror's Dryad culture metadata carries a
\texttt{DBTL\_Cycle} column, and Table~\ref{tab:pcsrawcycle} gives its distribution.
The join is on the construct name, so recovering the label is a preprocessing step
rather than a re-ingest. It is still a loader gap worth filing: a campaign whose entire
structure is the cycle should not lose the cycle on ingest.

\clearpage
%% GENERATED FILE -- do not hand-edit.
%% SOURCE: experiments/database/scripts/build_pputida_cell_state_table.py
\begingroup
\footnotesize
\setlength{\tabcolsep}{3pt}
\begin{longtable}{@{}L{44mm} L{62mm} L{62mm}@{}}
\caption[]{The two served isoprenol titer campaigns, axis by axis, every value read off the store. The product is the same compound in both, isoprenol at InChIKey CPJRRXSHAYUTGL-UHFFFAOYSA-N and PubChem 12988. \emph{Targets shared} is the intersection of the two CRISPRi target sets, counting the CRISPRi half of the genotype only, since a Carruthers record also carries the heterologous pathway.}\label{tab:pcsiso}\\
\toprule
\textbf{Axis} & \textbf{Carruthers 2025} & \textbf{Yunus 2026} \\
\midrule
\endfirsthead
\bottomrule
\endfoot
Records & 502 & 125 \\
\addlinespace[2pt]
CRISPRi targets & 135 & 121 \\
\addlinespace[2pt]
Targets shared by both campaigns & \multicolumn{2}{@{}l@{}}{120} \\
\addlinespace[2pt]
Genotype perturbation types & crispr interference, deletion strain, gene addition, heterologous pathway & crispr interference \\
\addlinespace[2pt]
Medium & M9-NREL + 20 g/L glucose (10 mM ammonium sulfate + Teknova T1001; Carruthers 2025) & M9 salts (6.78 g/L Na2HPO4, 3 g/L KH2PO4, 1 g/L NH4Cl, 0.5 g/L NaCl) \\
\addlinespace[2pt]
Temperature & 24\,\textdegree C & 30\,\textdegree C \\
\addlinespace[2pt]
Duration & 48 h & 48 h \\
\addlinespace[2pt]
Replicates & 2 biological\_replicate, 3 biological\_replicate, 6 biological\_replicate & 3 biological\_replicate \\
\addlinespace[2pt]
Quantification & GC-FID & GC-FID \\
\addlinespace[2pt]
Titer unit & ug/mL & ug/mL \\
\addlinespace[2pt]
Titer range & 0.00 to 771.2 & 1.59 to 1,495.0 \\
\addlinespace[2pt]
Records carrying an uncertainty & 502 of 502 & 0 of 125 \\
\addlinespace[2pt]
\end{longtable}
\endgroup
\noindent\emph{Targets in Yunus 2026 and not in Carruthers 2025:} \gene{PP\_\allowbreak 4129}.


%% GENERATED FILE -- do not hand-edit.
%% SOURCE: experiments/database/scripts/build_pputida_cell_state_table.py
\begingroup
\footnotesize
\begin{longtable}{@{}L{62mm} r@{\hspace{8pt}} r@{\hspace{8pt}} r@{}}
\caption[]{The DBTL cycle label in the served Carruthers stores. \emph{Naming a cycle} counts records whose full serialization matches \texttt{DBTL\textbackslash s*\textbackslash d} anywhere, so it is as permissive a search as the store allows.}\label{tab:pcscycle}\\
\toprule
\textbf{Store} & \textbf{Records} & \textbf{Naming a cycle} & \textbf{With a screen id} \\
\midrule
\endfirsthead
\bottomrule
\endfoot
campaign\_\allowbreak proteome\_\allowbreak carruthers2025 & 465 & 0 & 0 \\
isoprenol\_\allowbreak titer\_\allowbreak carruthers2025 & 502 & 0 & 0 \\
proteome\_\allowbreak carruthers2025 & 21 & 0 & 0 \\
proteome\_\allowbreak fold\_\allowbreak change\_\allowbreak carruthers2025 & 16 & 0 & 0 \\
\end{longtable}

\begingroup
\footnotesize
\begin{longtable}{@{}L{62mm} r@{}}
\caption[]{The same label in the raw mirror, which is the recovery path. Rows are cultures and include each cycle's controls, so these are not record counts; they establish that the label exists upstream and joins on \texttt{Line\_\allowbreak name}.}\label{tab:pcsrawcycle}\\
\toprule
\textbf{Cycle} & \textbf{Culture rows} \\
\midrule
\endfirsthead
\bottomrule
\endfoot
DBTL0 & 393 \\
DBTL1 & 183 \\
DBTL2 & 189 \\
DBTL3 & 174 \\
DBTL4 & 183 \\
DBTL5 & 186 \\
DBTL6 & 189 \\
\midrule
\textbf{Total} & \textbf{1,497} \\
\end{longtable}
\endgroup
\endgroup

